(1) 在远场观察方向，孔内横向位置 \(\boldsymbol r_\perp\) 的相位为 \(e^{i\boldsymbol q\cdot\boldsymbol r_\perp}\)，其中 \(q=k\sin\theta\)。利用极角积分 \(\int_0^{2\pi}e^{iqr\cos\phi}d\phi=2\pi J_0(qr)\)，远场复振幅正比
\[
2\pi\int_0^a rJ_0(qr)\,dr=\frac{2\pi aJ_1(qa)}q.
\]
令 \(u=qa=(2\pi a/\lambda)\sin\theta\)，以光轴强度归一化得到
\[
\frac{I(\theta)}{I(0)}
=\left[\frac{2J_1(u)}u\right]^2,
\qquad \lim_{u\to0}\frac{2J_1(u)}u=1.
\]

(2) 中央斑的第一个零强度圆环对应 \(J_1(u_1)=0\) 的首个正根 \(u_1\simeq3.83170597\)。因此
\[
\sin\theta_1=\frac{u_1\lambda}{2\pi a}
\simeq0.60983\frac{\lambda}{a}.
\]
在小角度、直径 \(D=2a\) 的情况下，\(\theta_1\simeq1.21967\,\lambda/D\)，俗称 Airy 衍射极限角半径。

(3) 若两个非相干点光源形成相同 Airy 强度分布，瑞利准则认为一个源的主峰恰落在另一源的第一暗环上时达到近似可分辨，故理想分辨角
\[
\Delta\theta_{\rm Rayleigh}\simeq1.22\frac{\lambda}{D}.
\]
这不是对任意信噪比、像差或检测算法都成立的绝对界限，而是经典衍射受限成像的约定准则。

(4) 在近轴远场中，环形能流的权重正比 \(u\,du\)。用 Bessel 恒等式 \(J_0'(u)=-J_1(u)\)、\(J_1'(u)=J_0(u)-J_1(u)/u\)，可验证
\[
\frac{d}{du}[1-J_0(u)^2-J_1(u)^2]=\frac{2J_1(u)^2}{u}.
\]
按所有角度归一化后累计强度正好为
\[
E(u)=1-J_0(u)^2-J_1(u)^2.
\]
当 \(u=u_1\)，因 \(J_1(u_1)=0\)，
\[
E(u_1)=1-J_0(u_1)^2\simeq0.8378 .
\]
约16.2%的光能位于第一个暗环外的衍射旁瓣，故中央斑虽然明亮并非包含全部能量。
